% Notes for the lawyer who reviews the FortyFive contracts. Not part of either agreement.
\documentclass[11pt]{article}
\usepackage{ff-contract}
\settitle{Notes for the lawyer}
\setversion{15 Oct 2026}
\fancyfoot[R]{\footnotesize Page \thepage{} of \pageref*{LastPage}}

\begin{document}

\titleblock{Notes for the lawyer}{What we decided, what we assumed, and what we need you to check in the Tutor Agreement and the Student and Parent Agreement}

\keynote{FortyFive is run by a 17-year-old founder, so the contracting entity is a parent's business (\ph{legal name}). Both agreements are plain-English first drafts, written for students, parents and young tutors to read in full. Please keep them readable when you change them.}

\section{Minors on both sides}
\begin{clauses}
  \item \textbf{Graduate tutors aged 17 (Tutor Agreement clause 4, Schedule 6).} The Guardian is the contracting party for money and liability; the minor signs only to acknowledge the rules; payouts go to the Guardian's account. Is this structure sound under section 11 of the Indian Contract Act, 1872? Does paid online tutoring by a 17-year-old raise anything under the Child and Adolescent Labour (Prohibition and Regulation) Act, 1986?
  \item \textbf{Students under 18 (Student and Parent Agreement).} The Parent is the party and pays; the Student signs the rules and the Coursework Code. Same question.
  \item \textbf{At 18.} Students choose whether a Parent keeps access (clause 4.5); graduate tutors re-sign on Track A (clause 4.6). Is a fresh agreement needed, or is ratification enough?
\end{clauses}

\section{Children's data}
\begin{clauses}
  \item Parental consent is planned through DigiLocker (Student and Parent Agreement clause 3.2). Does it meet ``verifiable consent'' under the Digital Personal Data Protection Act, 2023 and the DPDP Rules, 2025? Would a cheaper method also satisfy the Rules?
  \item Section 9(3) bars tracking, behavioural monitoring and targeted advertising of children. Do lesson recordings, the coursework Marker or the Platform's contact-detail filter in chat count as monitoring, and do they need an exemption or different wording?
  \item The 17-year-old tutor's own data is covered by the Guardian's consent (Schedule 6, item 2). Please confirm.
  \item Retention: recordings \ph{90 days}, Marker files 48 hours. Are these acceptable, and what must be kept longer for safety reports and tax?
\end{clauses}

\section{Safeguarding}
\begin{clauses}
  \item Our rule is to report suspected sexual offences straight to the police under sections 19 and 21 of the POCSO Act, 2012, with no internal inquiry first. Please check the wording in Tutor Agreement clause 9 and Student and Parent Agreement clause 11.
  \item Immediate suspension of a tutor with payouts held (Tutor Agreement clauses 9.5, 14.3). Is this fair and enforceable?
  \item Track B tutors have a school reference instead of police verification. Can a 17-year-old get police verification at all, and should insurance cover change for this track?
\end{clauses}

\section{Tutors as independent contractors}
\begin{clauses}
  \item Are our tutors ``platform workers'', and is FortyFive an ``aggregator'', under the Code on Social Security, 2020? If so, what contribution or registration applies, and from when?
  \item Clause 13.2 bars contact with Under-18s met through FortyFive for 12 months after the agreement ends. We framed it as a safeguarding rule, but section 27 of the Indian Contract Act, 1872 may make it unenforceable as a restraint of trade. Should it stay, change, or move into the Code of Conduct as a conduct rule?
  \item Track C teachers may not take paid sessions with students they teach or assess at their own school (clause 5.7). Is this enough to manage the conflict of interest?
\end{clauses}

\section{Tax and payments}
\begin{clauses}
  \item TDS: we wrote 0.1\% under the Income-tax Act, 2025 provision that replaced section 194-O of the 1961 Act, with the individual exemption below Rs 5 lakh for tutors who give a PAN (Tutor Agreement clause 6.7). Please confirm the section number and rate.
  \item GST: FortyFive as an e-commerce operator, GST on our fee and on Course sales, TCS under section 52 of the CGST Act for registered tutors. Is the receipt in Schedule 2 correct, and should it show the tutor's GST status?
  \item Founding tutors pay no fee until 30 April 2027 (clause 6.4). Any GST or valuation issue with a zero fee?
\end{clauses}

\section{Consumers, refunds and disputes}
\begin{clauses}
  \item Grievance Officer, the 48-hour acknowledgement and the one-month resolution under the Consumer Protection (E-Commerce) Rules, 2020. What else must we display?
  \item Refunds: half the Price for a cancellation within 24 hours, and a 15-minute no-show rule. These are business choices; are they fair under the Consumer Protection Act, 2019?
  \item The tutor agreement uses arbitration seated in Coimbatore. The student agreement keeps the consumer forums open instead. Is that the right split?
  \item Our liability caps: 12 months of fees (tutors), 3 months of payments (families). Are they enforceable?
\end{clauses}

\section{Content and intellectual property}
\begin{clauses}
  \item We never use IB papers, markschemes or criteria text, and course leads write in their own words. Is our use of ``IB'' in descriptions, and the disclaimer, enough to keep clear of the IBO's trademarks?
  \item Course leads give a \ph{2-year} exclusive licence (Tutor Agreement Schedule 4, item 5). Is the length reasonable?
\end{clauses}

\section{Signing}
\begin{clauses}
  \item We plan click-to-sign inside the sign-up flow, with Aadhaar eSign as an option. Which do you recommend for each agreement, and does Tamil Nadu stamp duty apply?
  \item Who in the parent's business signs for FortyFive, and what authority do they need?
\end{clauses}

\end{document}
